\paragraph{Data.} Classification uses five public CT cohorts with \numLabelsN{} binary labels (Table~\ref{tab:cohorts}): the CVPR26 cohort of the CVPR 2026 3-D CT foundation-model challenge, which groups AMOS \citep{ji2022amos} with 15 disease labels, COVID-CTset \citep{rahimzadeh2021covidctset} and LUNA25 nodule malignancy \citep{peeters2026luna25}; FDG-PET/CT \citep{gatidis2022fdgpetct}, whose four mutually exclusive diagnosis groups we predict from the CT alone as one-vs-rest labels; RSNA 2022 cervical-spine fractures \citep{lin2023rsnacsfdata} (patient level and C1--C7); RSNA 2023 abdominal trauma \citep{rudie2024ratic} (any injury and kidney, liver and spleen injury); and RSNA 2019 intracranial hemorrhage \citep{flanders2020rsnaich} (any hemorrhage and five subtypes). Segmentation uses AMOS22 (15 organs; \citealp{ji2022amos}), KiTS23\footnote{\url{https://github.com/neheller/kits23}} (kidney, tumor, cyst; \citealp{heller2023kits21}), LiTS (liver, tumor; \citealp{bilic2023lits}), MSD Pancreas (pancreas, tumor; \citealp{antonelli2022msd}) and TotalSegmentator \citep{wasserthal2023totalsegmentator,wasserthal2026totalsegdata}, whose 117 structures we merge within anatomical systems into 15 (Appendix~\ref{app:data}); \numStructuresN{} structures in all. Every source is divided into a training, a selection and a held-out split: for the CVPR26 cohort and AMOS22 the official validation cases are divided in halves, for the other sources the splits are fixed seeded divisions of the released labelled cases, with the official TotalSegmentator test split held out. Patients do not span splits where patient identifiers exist (FDG-PET/CT, RSNA 2023, RSNA 2019); the other sources are split per case. Pillar-0 is run on the hemorrhage cohort, and VoCo on FDG-PET/CT, RSNA 2022 and RSNA 2023 and on all segmentation datasets.

\paragraph{Training.} Every run uses AdamW with a linear warm-up over the first 10\% of the steps, cosine decay to $10^{-6}$, bf16 mixed precision and seed 42. A classification run trains 8000 steps at batch size 1 on a masked, positive-weighted binary cross-entropy; a segmentation run trains 8000 steps of four crops on soft Dice plus binary cross-entropy per structure; two-stage methods split the budget into halves. The joint models of the two 15-structure datasets (AMOS22, TotalSegmentator) train 24\,000 steps, and the CVPR26 joint model is additionally trained for 26\,000 steps as a longer schedule reported in the appendix. Each method has one default learning rate and weight decay, set once before the sweep by the type of parameter it trains, from the largest rate for a head on frozen features ($10^{-3}$; $10^{-2}$ for the segmentation linear probe) through low-rank adapters ($10^{-4}$) to the smallest for encoder weights ($2\times10^{-5}$), and not tuned on our data for any backbone, cohort or method (Tables~\ref{tab:optim} and~\ref{tab:optimseg}). Every 1000 steps both the selection and the held-out split are evaluated; the checkpoint with the best selection macro AUROC or mean Dice is kept and its held-out metric is reported. TotalSegmentator per-structure models use Lin, Dec, LoRA and Full FT only. The joint models of Lin, LoRA, Attn-only FT and Full FT (classification) and of Dec, LoRA and Full FT (segmentation) were trained with two further seeds on TAP-CT, SPECTRE and Pillar-0 (\numSeedsclsRunsN{} and \numSeedssegRunsN{} runs). All settings are listed in Appendix~\ref{app:details}.

\paragraph{Evaluation and statistics.} Held-out AUROC per label and Dice per structure and case are the metrics \citep{maierhein2024metrics}. Per label--mode pair, methods are compared with DeLong's test on AUROC \citep{delong1988} and with a paired case bootstrap on Dice, with Benjamini--Hochberg correction \citep{benjamini1995fdr} within each family of comparisons at $\alpha=0.05$, and ranked within each pair that has every method. The main comparisons are summarised at label level: each label counts once with its per-task and joint differences averaged, and the 95\% interval of the mean comes from a percentile bootstrap over cohorts, labels within cohorts and held-out cases within labels (Appendix~\ref{app:stats}); because four or five cohorts per backbone make the cohort stage optimistic, we also report the number of cohorts with a positive mean difference. Pooled estimates use TAP-CT, SPECTRE and Pillar-0; VoCo, which ran three of the five classification cohorts, is reported per backbone. A method whose held-out probabilities barely vary between cases did not train: we flag it when the median standard deviation of its held-out probabilities over label--mode pairs is below 0.01, keep it in the per-cohort tables and leave it out of ranks and averages over methods.
